% !TeX root = ../../Book2.tex
\chapter{热传导方程}
\label{b2:ch:12}
\BookChapterNumber{b2:ch:12}
% 本章摘要（不计入编号）
\begin{chapterabstract}
热传导把空间振荡逐渐消去。有限区间上，边界条件决定离散模态及最终平衡；全空间上，Gauss 核把扩散写成平移平均。本章将前面的特征展开与 Fourier 变换用于初边值问题，说明粗初温怎样在正时间变为光滑函数，并分别在有界经典解类与平方可积解类中证明唯一性。最后把时间演化组织成热半群，明确收缩性、强连续性的空间范围，以及生成元的定义域。
\end{chapterabstract}

% 本章目标（不计入编号）
\begin{learninggoals}
\begin{itemize}
\item\label{b2:item:12-goal-boundary}用特征展开与镜像核求解区间热问题，区分固定端温、绝热端点和散热边界。
\item\label{b2:item:12-goal-kernel}从已知 Gauss 变换取得全空间核解，并证明初值收敛和含源项公式。
\item\label{b2:item:12-goal-smooth}按抛物尺度估计各阶导数，解释光滑化与无穷传播速度。
\item\label{b2:item:12-goal-unique}使用最大值原理或平方能量控制唯一性，说明全空间的解类限制。
\item\label{b2:item:12-goal-semigroup}检查热半群的时间连续性，识别 \(L^2\) 生成元及反向时间的不稳定性。
\end{itemize}
\end{learninggoals}

区间部分直接使用\chapref{b2:ch:10}的离散谱、能量空间和热级数极限；全空间部分使用\chapref{b2:ch:03}的 Gauss 变换与逼近恒等、\chapref{b2:ch:04}的 Plancherel 理论。\chapref{b2:ch:06}的圆周周期化提供镜像核。\chapref{b2:ch:11}的传播现象可以作比较，但不参与本章存在性证明。

\ProseParagraphBreak

\secref{b2:sec:1201}先从边界与长期平衡理解扩散，\secref{b2:sec:1202}给出全空间表示，\secref{b2:sec:1203}分析正则性、比较与唯一性，\secref{b2:sec:1204}统一时间观点。Lawrence C. Evans 的《Partial Differential Equations》\cite{Evans2010Partial}可作为热方程、能量方法与线性发展方程的系统伴读；Walter A. Strauss 的《Partial Differential Equations: An Introduction》\cite{Strauss2008Partial}适合结合有限杆模型练习边界处理。此处只推荐两书的主题范围，不以未取得的正文页码替代下文论证。

% !TeX root = ../../Book2.tex
\section{分离变量法}
\label{b2:sec:1201}

\subsection{有界区间}
\label{b2:subsec:120101}

热传导中，温度梯度引起从高温向低温的通量。设细杆的热容量密度为 \(w(x)>0\)、传导系数为 \(p(x)>0\)，通量取 \(-p(x)u_x\)。无热源时，对任意子区间的热量平衡给出 \(wu_t=(pu_x)_x\)。若再有与温度成正比的损耗 \(q(x)u\)，则方程成为
\[
 wu_t=(pu_x)_x-qu,
 \quad u_t+Au=0.
\]
在有限区间及\cref{b2:def:10-sl-operator}的连续系数、分离边界假设下，空间算子就是\chapref{b2:ch:10}建立的 Sturm--Liouville 算子。本节先取非负能量情形；例如 \(q\geq0\)，以及 Robin 系数非负。

\ProseParagraphBreak

对初温 \(f\in H=L^2(I,w\,\mathrm{d}x)\)，\cref{b2:prop:10-heat-energy-limit}给出
\begin{equation}\label{b2:eq:12-interval-expansion}
 u(t)=\sum_ne^{-\lambda_nt}f_ne_n,
 \quad f_n=\langle f,e_n\rangle_H.
\end{equation}
它在零时以 \(H\) 范数取得初值，正时间属于 \(D(A)\)，并在 \(H\) 中满足方程；还在 \(L^2((0,T);V)\) 中具有有限能量。连续系数并不使空间解自动无限次可微：当 \(p\) 只有连续性时，正则性应以 \(pu_x\) 和 \(D(A)\) 表述。下面转到常系数杆，才进一步证明所有空间阶数的经典光滑性。

\subsection{Fourier 级数}
\label{b2:subsec:120102}

在 \(0<x<\ell\) 上固定 \(\kappa>0\)，考虑 \(u_t=\kappa u_{xx}\)。齐次 Dirichlet 条件给出
\begin{equation}\label{b2:eq:12-dirichlet-series}
 u_D(t,x)=\sum_{n=1}^infty f_ne^{-\kappa(n\pi/\ell)^2t}
 \sqrt{\dfrac2\ell}\sin\dfrac{n\pi x}{\ell}.
\end{equation}
齐次 Neumann 条件给出
\begin{equation}\label{b2:eq:12-neumann-series}
 u_N(t,x)=\overline f+
 \sum_{n=1}^infty g_ne^{-\kappa(n\pi/\ell)^2t}
 \sqrt{\dfrac2\ell}\cos\dfrac{n\pi x}{\ell},
 \quad\overline f=\dfrac1\ell\int_0^\ell f.
\end{equation}
其中 \(f_n,g_n\) 分别为正弦和非零余弦系数。常数项来自零模态；它不随时间衰减。

\begin{proposition}[正时间的经典解]\label{b2:prop:12-interval-smooth}
若 \(f\in L^2(0,\ell)\)，上面两个级数及其任意有限阶 \(x,t\) 导数在 \([\delta,T]\times[0,\ell]\) 上绝对一致收敛，其中 \(0<\delta<T\)。因此它们在正时间光滑，满足各自方程和边界条件。
\end{proposition}
\begin{proof}
求 \(j\) 次空间导数及 \(m\) 次时间导数只会把第 \(n\) 项乘以至多 \(C_{j,m}n^{j+2m}\)。Cauchy--Schwarz 给出统一可和上界
\[
 \sum_{n\geq1}|f_n|n^{j+2m}e^{-c\delta n^2}
 \leq\left(\sum_{n\geq1}|f_n|^2\right)^{1/2}
 \left(\sum_{n\geq1}n^{2j+4m}e^{-2c\delta n^2}\right)^{1/2}<\infty.
\]
余弦系数的证明相同。每个阶数都有相应一致收敛，故逐项微分合法。有限项已经满足方程和边界条件，极限保持它们。
\end{proof}

例如两端温度固定为零，初温在杆内恒为 \(1\)，则
\[
 u_D(t,x)=\dfrac4\pi\sum_{j=0}^{\infty}
 \dfrac{e^{-\kappa((2j+1)\pi/\ell)^2t}}{2j+1}
 \sin\dfrac{(2j+1)\pi x}{\ell}.
\]
正时间的端点值为零，而零时的 \(L^2\) 初温仍是常数一。它不可能在闭区间上随着 \(t\downarrow0\) 一致趋于一，因为端点值始终为零；初边值在角点不相容并不妨碍平方能量解存在。

\subsection{边界条件}
\label{b2:subsec:120103}

周期热核提供另一种表示。沿用第一编的实线核 \(H_s(x)=(4\pi s)^{-1/2}e^{-x^2/(4s)}\)，定义
\begin{align*}
 K_D(t,x,y)&=\sum_{m\in\mathbb Z}
 \left(H_{\kappa t}(x-y+2m\ell)-H_{\kappa t}(x+y+2m\ell)\right),\\
 K_N(t,x,y)&=\sum_{m\in\mathbb Z}
 \left(H_{\kappa t}(x-y+2m\ell)+H_{\kappa t}(x+y+2m\ell)\right).
\end{align*}
则 \(u_D(t,x)=\int_0^\ell K_D(t,x,y)\cdot f(y)\,\mathrm{d}y\)，Neumann 情形用 \(K_N\)。两式分别来自 \(f\) 的奇延拓、偶延拓及 \(2\ell\) 周期化；将卷积在每个周期分割后换元即得。高斯尾保证对正时间紧区间的一切求和与微分可交换。第一编的周期 Fourier 系数唯一性证明它们与式\eqref{b2:eq:12-dirichlet-series}、式\eqref{b2:eq:12-neumann-series}一致。

\ProseParagraphBreak

这两个空间表达式在小时间有直接用途：附近的镜像点控制边界影响，远处项指数小。大时间则应使用特征级数，只保留低频项。由于奇延拓与偶延拓在有限 \(p\) 范数中有界，周期热核逼近还给出 \(f\in L^p(0,\ell)\)、\(1\leq p<\infty\) 时的初值收敛。Dirichlet 的闭区间一致收敛要求 \(f\) 连续且两端为零；Neumann 情形对任意连续 \(f\) 都成立，零时不要求导数边界相容。

\ProseParagraphBreak

非齐次端点先用提升处理。例如 \(u(t,0)=a(t)\)、\(u(t,\ell)=b(t)\)，其中 \(a,b\in C^1\)，令
\[
 r(t,x)=a(t)\left(1-\dfrac x\ell\right)+b(t)\dfrac x\ell,
 \quad v=u-r.
\]
那么 \(v\) 的边界为零，且 \(v_t-\kappa v_{xx}=-r_t\)。记源项正弦系数为 \(F_n(t)\)，逐模态公式为
\[
 v_n(t)=e^{-\kappa\lambda_nt}v_n(0)
 +\int_0^te^{-\kappa\lambda_n(t-s)}F_n(s)\,\mathrm{d}s.
\]
对 \(F\in C([0,T];L^2)\)，积分项的平方范数由 \(\int_0^t\|F(s)\|_2\,\mathrm{d}s\) 控制，有限截断因紧时间区间上的连续像为紧集而一致逼近。因而该公式给出 \(C([0,T];L^2)\) 中的解；不能只因源项连续就额外断言所有高阶导数连续。

\ProseParagraphBreak

Robin 条件的解仍用式\eqref{b2:eq:12-interval-expansion}。热量是否守恒取决于通量：常系数时
\[
 \dfrac{\mathrm{d}}{\mathrm{d}t}\int_0^\ell u
 =\kappa\left(u_x(t,\ell)-u_x(t,0)\right).
\]
Neumann 的两项为零；Robin 的两项给出端点散热；Dirichlet 固定端温却不固定通量，不能由齐次边界推出总热量不变。

\subsection{长时间行为}
\label{b2:subsec:120104}

\begin{proposition}[趋向零模态]\label{b2:prop:12-long-time}
设本节的空间算子非负，\(P_0\) 为到零特征空间的正交投影；若没有零特征值则 \(P_0=0\)。记最小正特征值为 \(\lambda_+\)。则
\[
 \|u(t)-P_0f\|_H\leq e^{-\lambda_+t}\|f-P_0f\|_H.
\]
常系数 Dirichlet 情形 \(P_0f=0\)，Neumann 情形 \(P_0f=\overline f\)；两者 \(\lambda_+=\kappa\pi^2/\ell^2\)。
\end{proposition}
\begin{proof}
将零模态从展开中分离，剩余每个系数的平方都乘以不大于 \(e^{-2\lambda_+t}\) 的因子，求和即得。最小正特征值存在由\cref{b2:thm:10-sl-spectral}保证。
\end{proof}

对常系数区间，这还给出 \(t\geq1\) 时的一致收敛：先演化至 \(t-1\)，再由固定时间一的核把 \(L^2\) 误差送到一致范数；Cauchy--Schwarz 与正时间可和尾给出固定的 \(L^2\to L^\infty\) 界。若初值恰与最低正模态正交，衰减速度由下一项非零系数决定。正 Robin 系数排除常数零模态，故最终冷却到零；出现负特征值时却可有指数增长，不能继续使用本命题。


% !TeX root = ../../Book2.tex
\section{Fourier 变换法}
\label{b2:sec:1202}

\subsection{热核}
\label{b2:subsec:120201}

全空间没有端点，也没有正的最低允许频率。Fourier 变换把每个连续频率单独演化。先对足够光滑且衰减的解形式计算，按照本卷的 \(2\pi\) 约定有
\[
 \partial_t\widehat u(t,\xi)
 =-4\pi^2\kappa|\xi|^2\widehat u(t,\xi),
 \quad\widehat u(t,\xi)=e^{-4\pi^2\kappa t|\xi|^2}\widehat f(\xi).
\]
\cref{b2:thm:03-gaussian}已经计算了这个乘子的逆变换。沿用第一编的热核记号，在 \(\mathbb R^d\) 上写作
\begin{equation}\label{b2:eq:12-gauss-heat}
 H_t(x)=(4\pi t)^{-d/2}e^{-|x|^2/(4t)},
 \quad t>0,
 \quad\widehat H_t(\xi)=e^{-4\pi^2t|\xi|^2}.
\end{equation}
方程扩散系数为 \(\kappa\) 时使用 \(H_{\kappa t}\)，不另换 Fourier 归一化。这个核非负、总积分为一；其空间宽度约为 \(\sqrt t\)，高度约为 \(t^{-d/2}\)。由直接微分有 \(\partial_tH_{\kappa t}=\kappa\Delta H_{\kappa t}\)。

\ProseParagraphBreak

这些计算可以独立验证卷积候选，不要求先假定未知解的变换存在。对于粗初值，求解的逻辑次序是先用核定义 \(u\)，再证明可微性和初值收敛；前面的频率微分只负责找到候选。第一编已经证明 Gauss 核的逼近恒等性质，下面把其参数取为 \(\sqrt{4\pi\kappa t}\)。

\subsection{卷积表示}
\label{b2:subsec:120202}

\begin{theorem}[全空间热核解]\label{b2:thm:12-heat-convolution}
设 \(1\leq p\leq\infty\)、\(f\in L^p(\mathbb R^d)\)。则
\begin{equation}\label{b2:eq:12-heat-convolution}
 u(t,x)=\int_{\mathbb R^d}H_{\kappa t}(x-y)\cdot f(y)\,\mathrm{d}y
\end{equation}
对 \(t>0\)、\(x\in\mathbb R^d\) 有意义，且在 \((0,\infty)\times\mathbb R^d\) 上光滑，满足 \(u_t=\kappa\Delta u\)。当 \(p<\infty\) 时，\(u(t)\to f\) 于 \(L^p\)；当 \(f\) 有界一致连续时，收敛在一致范数中成立。若有界函数 \(f\) 在 \(x_0\) 连续，则 \(u(t,x_0)\to f(x_0)\)。
\end{theorem}
\begin{proof}
热核及其每阶时空导数属于每个 \(L^r\)、\(1\leq r\leq\infty\)。固定 \(t\) 的正紧区间与 \(x\) 的紧集后，每个导数可由一个平移后仍快速衰减的函数统一控制。对 \(f\in L^p\)，Hölder 不等式给出绝对可积性；对差商使用同一控制，可以把每一阶导数移入积分。因此卷积光滑，且核的微分等式传给 \(u\)。这里在 \(p=1\) 用核导数的一致界，在 \(p=\infty\) 用其绝对可积性，两个端点都被覆盖。

有限 \(p\) 与一致连续情形由\cref{b2:thm:03-approximate-identity}得到。对指定连续点，把 \(f(x_0-y)-f(x_0)\) 的积分分成 \(|y|<\delta\) 与补集。近处由连续性控制，远处由 \(2\|f\|_\infty\) 乘热核尾质量控制，先选 \(\delta\) 再令 \(t\downarrow0\) 即得。
\end{proof}

初值也可以是一点质量，此时解是 \(H_{\kappa t}(x-y_0)\)，在 \(t\downarrow0\) 时只以测度或分布意义趋于 \(\delta_{y_0}\)。为了区别空间核与时空基本解，令 \(E(t,x)=H_{\kappa t}(x)\) 当 \(t>0\)，在 \(t<0\) 令 \(E=0\)。那么
\[
 (\partial_t-\kappa\Delta_x)E=\delta_{(0,0)}
\]
在时空分布中成立。确实，对紧支集光滑 \(\varphi(t,x)\)，从 \(t=\varepsilon\) 起分部积分，热方程抵消内部项，留下 \(\int H_{\kappa\varepsilon}(x)\cdot\varphi(\varepsilon,x)\,\mathrm{d}x\)。它由逼近恒等趋于 \(\varphi(0,0)\)。同时 \(E\) 在有限时间带的空间积分为该带长度，所以零时附近的局部可积性也已得到。

\ProseParagraphBreak

对非齐次方程，设 \(1\leq p<\infty\)、\(F\in C([0,T];L^p)\)。由逐时叠加得到
\begin{equation}\label{b2:eq:12-duhamel}
 u(t)=H_{\kappa t}*f+
 \int_0^tH_{\kappa(t-s)}*F(s)\,\mathrm{d}s.
\end{equation}
在积分的右端按 \(H_0*F(t)=F(t)\) 理解。被积函数是连续的 \(L^p\) 值函数，积分可作为 Riemann 和的范数极限；Young 不等式给出
\[
 \left\|\int_0^tH_{\kappa(t-s)}*F(s)\,\mathrm{d}s\right\|_p
 \leq\int_0^t\|F(s)\|_p\,\mathrm{d}s.
\]
这保证源项趋于零初值。先对光滑紧支集的时空源项求导，再以连续 \(L^p\) 范数逼近一般 \(F\)，便在分布中得到 \(u_t-\kappa\Delta u=F\)。逼近可将时间分成有限小段，对有限个空间取值分别磨光，再以时间连续插值连接；上式使解同时收敛，测试函数的积分配对也连续。这个验证不把一般连续源项的积分公式误写成已经证明的经典二阶正则性。

\subsection{高维热方程}
\label{b2:subsec:120203}

高维热核是各坐标一维核的乘积，但高维扩散的幅度不能从一维直观中省略维数。换元 \(x=\sqrt t\,z\) 给出
\[
 \|H_t\|_r=C_{d,r}t^{-\frac d2(1-1/r)},
 \quad1\leq r\leq\infty.
\]
例如有限质量初温满足 \(\|u(t)\|_\infty\leq(4\pi\kappa t)^{-d/2}\|f\|_1\)。没有任何质量损失也能使温度峰值降低，因为同一热量散布到越来越大的区域。

\ProseParagraphBreak

方向上的扩散速率也可以不同。若 \(B\) 是实对称正定常矩阵，方程为 \(u_t=\sum_{i,j}B_{ij}\partial_i\partial_ju\)，则相应核是
\[
 H_t^B(x)=\dfrac1{(4\pi t)^{d/2}\sqrt{\det B}}
 \exp\left(-\dfrac{x^{\mathsf T}B^{-1}x}{4t}\right).
\]
正交对角化 \(B\) 后，各方向分别使用一维扩散核，换元的行列式给出归一化因子。故核仍非负、质量为一，变换为 \(e^{-4\pi^2t\xi^{\mathsf T}B\xi}\)。正定性是每个方向都得到磨光的条件；若某个特征值为零，该方向不发生卷积扩散，不能继续声称所有空间方向光滑。

\ProseParagraphBreak

这一正定二次型公式可与 John K. Hunter、Bruno Nachtergaele 的《Applied Analysis》\cite[pp. 303--304]{Hunter2001Applied}对照，其变换采用不同的 \(2\pi\) 归一化。本卷已经在\secref{b2:sec:0502}换算过常数，这里直接使用一致的频率约定。

% !TeX root = ../../Book2.tex
\section{热方程的结构}
\label{b2:sec:1203}

显式热核已经给出候选解。现在要分别检查正时间的导数、最大值控制和解的唯一性，并把光滑计算延伸到平方可积初值及时间可积的外力。每一种结论所用的解类都会保留在陈述中，尤其是全空间无穷远处的控制。

\subsection{光滑化效应}
\label{b2:subsec:120301}

热核的任意阶导数仍是多项式乘 Gauss 函数，因此正时间磨光不仅是定性结论。一次空间求导损失半个时间尺度，一次时间求导损失一个时间尺度；这是 \(x\sim\sqrt t\) 的抛物缩放直接产生的。

\Needspace{6\baselineskip}
\begin{theorem}[热流的导数估计]\label{b2:thm:12-smoothing}
设整数 \(d\geq1\)、\(\kappa>0\)，\(H_s(x)=(4\pi s)^{-d/2}e^{-|x|^2/(4s)}\)（\(s>0\)、\(x\in\mathbb R^d\)）。设 \(f\in L^p(\mathbb R^d)\)，\(1\leq p\leq q\leq\infty\)，\(\alpha\) 为多重指标，\(m\geq0\) 为整数。对 \(u(t)=H_{\kappa t}*f\)，有
\begin{equation}\label{b2:eq:12-smoothing}
 \|\partial_t^m\partial_x^\alpha u(t)\|_q
 \leq C t^{-m-|\alpha|/2-\frac d2(1/p-1/q)}\|f\|_p,
 \quad t>0.
\end{equation}
常数只依赖所列指数、维数与 \(\kappa\)，不依赖 \(t,f\)。
\end{theorem}
\begin{proof}
缩放与直接求导给出
\[
 \partial_t^m\partial_x^\alpha H_{\kappa t}(x)
 =t^{-m-|\alpha|/2-d/2}\Psi_{m,\alpha}(x/\sqrt t),
\]
其中 \(\Psi_{m,\alpha}\) 是快速衰减的光滑函数。因此其 \(L^r\) 范数为常数乘 \(t^{-m-|\alpha|/2-\frac d2(1-1/r)}\)。Young 不等式给出 \(p\to p\) 的导数界，逐点 Hölder 不等式给出 \(p\to\infty\) 的界。若 \(p<\infty\)、\(p\leq q<\infty\)，对导数函数 \(v\) 使用
\[
 \|v\|_q\leq\|v\|_p^{p/q}\cdot\|v\|_\infty^{1-p/q},
\]
由两个端点界即得所述指数。\(p=\infty\) 时只能取 \(q=\infty\)，已含在 Young 估计中。
\end{proof}

当 \(t\downarrow0\) 时，导数界可以发散，这是粗初值与正时间光滑性共存的方式。若初值已有某个 \(L^p\) 弱导数，可把该导数移到 \(f\) 上，就不必支付相应的时间奇性。取一维跳跃初温 \(f=\mathbf1_{(0,\infty)}\)，解为
\[
 u(t,x)=\dfrac1{\sqrt\pi}\int_{-\infty}^{x/\sqrt{4\kappa t}}e^{-s^2}\,\mathrm{d}s,
 \quad u_x(t,x)=H_{\kappa t}(x).
\]
跃变被摊开到 \(\sqrt{\kappa t}\) 的过渡层，梯度峰值为 \((4\pi\kappa t)^{-1/2}\)。这同时给出不能在整个 \(L^\infty\) 中要求初值一致收敛的具体例子。\Cref{b2:fig:12-heat-step}将三个正时间的剖面与初温并列：曲线随时间变缓，过渡层随之变宽；在初温为零的一侧，正时间已处处出现非零的细小尾部。
% Step heat profiles: values sampled directly from the proved Gaussian integral, kappa=1.
\begin{figure}[htbp]
\centering
\begin{tikzpicture}[x=1.1cm,y=3.4cm]
\draw[->,AnalysisGray] (-3.35,0)--(3.35,0) node[right]{\(x\)};
\draw[->,AnalysisGray] (-3.3,0)--(-3.3,1.15) node[above]{\(u\)};
\foreach \x in {-3,-2,-1,0,1,2,3}{
 \draw (\x,0)--(\x,-.02) node[below,font=\small]{\(\x\)};
}
\foreach \y/\lab in {.25/{\frac14},.5/{\frac12},.75/{\frac34},1/1}{
 \draw (-3.3,\y)--(-3.36,\y) node[left,font=\small]{\(\lab\)};
}
\draw[AnalysisGray!35,thin,densely dotted] (-3.28,.5)--(3.15,.5);
\draw[black,thick,densely dotted] (-3,0)--(0,0);
\draw[black,thick,densely dotted] (0,1)--(3,1);
\fill (0,0) circle(1.3pt);
\draw[fill=white] (0,1) circle(1.3pt);
\draw[AnalysisBlue,very thick] plot coordinates {
 (-3.0000,0.00000000) (-2.9700,0.00000000) (-2.9400,0.00000000) (-2.9100,0.00000000) (-2.8800,0.00000000) (-2.8500,0.00000000) (-2.8200,0.00000000) (-2.7900,0.00000000)
 (-2.7600,0.00000000) (-2.7300,0.00000000) (-2.7000,0.00000000) (-2.6700,0.00000000) (-2.6400,0.00000000) (-2.6100,0.00000000) (-2.5800,0.00000000) (-2.5500,0.00000000)
 (-2.5200,0.00000000) (-2.4900,0.00000000) (-2.4600,0.00000000) (-2.4300,0.00000000) (-2.4000,0.00000000) (-2.3700,0.00000000) (-2.3400,0.00000000) (-2.3100,0.00000000)
 (-2.2800,0.00000000) (-2.2500,0.00000000) (-2.2200,0.00000000) (-2.1900,0.00000000) (-2.1600,0.00000000) (-2.1300,0.00000000) (-2.1000,0.00000000) (-2.0700,0.00000000)
 (-2.0400,0.00000000) (-2.0100,0.00000001) (-1.9800,0.00000001) (-1.9500,0.00000002) (-1.9200,0.00000003) (-1.8900,0.00000005) (-1.8600,0.00000007) (-1.8300,0.00000011)
 (-1.8000,0.00000018) (-1.7700,0.00000028) (-1.7400,0.00000043) (-1.7100,0.00000066) (-1.6800,0.00000101) (-1.6500,0.00000153) (-1.6200,0.00000230) (-1.5900,0.00000344)
 (-1.5600,0.00000511) (-1.5300,0.00000754) (-1.5000,0.00001105) (-1.4700,0.00001607) (-1.4400,0.00002321) (-1.4100,0.00003330) (-1.3800,0.00004746) (-1.3500,0.00006717)
 (-1.3200,0.00009441) (-1.2900,0.00013180) (-1.2600,0.00018275) (-1.2300,0.00025168) (-1.2000,0.00034426) (-1.1700,0.00046772) (-1.1400,0.00063117) (-1.1100,0.00084603)
 (-1.0800,0.00112642) (-1.0500,0.00148973) (-1.0200,0.00195710) (-0.9900,0.00255400) (-0.9600,0.00331089) (-0.9300,0.00426376) (-0.9000,0.00545475) (-0.8700,0.00693270)
 (-0.8400,0.00875361) (-0.8100,0.01098096) (-0.7800,0.01368594) (-0.7500,0.01694743) (-0.7200,0.02085172) (-0.6900,0.02549198) (-0.6600,0.03096742) (-0.6300,0.03738203)
 (-0.6000,0.04484301) (-0.5700,0.05345884) (-0.5400,0.06333692) (-0.5100,0.07458099) (-0.4800,0.08728818) (-0.4500,0.10154589) (-0.4200,0.11742864) (-0.3900,0.13499478)
 (-0.3600,0.15428344) (-0.3300,0.17531166) (-0.3000,0.19807195) (-0.2700,0.22253037) (-0.2400,0.24862516) (-0.2100,0.27626619) (-0.1800,0.30533515) (-0.1500,0.33568662)
 (-0.1200,0.36714997) (-0.0900,0.39953208) (-0.0600,0.43262082) (-0.0300,0.46618920) (0.0000,0.50000000) (0.0300,0.53381080) (0.0600,0.56737918) (0.0900,0.60046792)
 (0.1200,0.63285003) (0.1500,0.66431338) (0.1800,0.69466485) (0.2100,0.72373381) (0.2400,0.75137484) (0.2700,0.77746963) (0.3000,0.80192805) (0.3300,0.82468834)
 (0.3600,0.84571656) (0.3900,0.86500522) (0.4200,0.88257136) (0.4500,0.89845411) (0.4800,0.91271182) (0.5100,0.92541901) (0.5400,0.93666308) (0.5700,0.94654116)
 (0.6000,0.95515699) (0.6300,0.96261797) (0.6600,0.96903258) (0.6900,0.97450802) (0.7200,0.97914828) (0.7500,0.98305257) (0.7800,0.98631406) (0.8100,0.98901904)
 (0.8400,0.99124639) (0.8700,0.99306730) (0.9000,0.99454525) (0.9300,0.99573624) (0.9600,0.99668911) (0.9900,0.99744600) (1.0200,0.99804290) (1.0500,0.99851027)
 (1.0800,0.99887358) (1.1100,0.99915397) (1.1400,0.99936883) (1.1700,0.99953228) (1.2000,0.99965574) (1.2300,0.99974832) (1.2600,0.99981725) (1.2900,0.99986820)
 (1.3200,0.99990559) (1.3500,0.99993283) (1.3800,0.99995254) (1.4100,0.99996670) (1.4400,0.99997679) (1.4700,0.99998393) (1.5000,0.99998895) (1.5300,0.99999246)
 (1.5600,0.99999489) (1.5900,0.99999656) (1.6200,0.99999770) (1.6500,0.99999847) (1.6800,0.99999899) (1.7100,0.99999934) (1.7400,0.99999957) (1.7700,0.99999972)
 (1.8000,0.99999982) (1.8300,0.99999989) (1.8600,0.99999993) (1.8900,0.99999995) (1.9200,0.99999997) (1.9500,0.99999998) (1.9800,0.99999999) (2.0100,0.99999999)
 (2.0400,1.00000000) (2.0700,1.00000000) (2.1000,1.00000000) (2.1300,1.00000000) (2.1600,1.00000000) (2.1900,1.00000000) (2.2200,1.00000000) (2.2500,1.00000000)
 (2.2800,1.00000000) (2.3100,1.00000000) (2.3400,1.00000000) (2.3700,1.00000000) (2.4000,1.00000000) (2.4300,1.00000000) (2.4600,1.00000000) (2.4900,1.00000000)
 (2.5200,1.00000000) (2.5500,1.00000000) (2.5800,1.00000000) (2.6100,1.00000000) (2.6400,1.00000000) (2.6700,1.00000000) (2.7000,1.00000000) (2.7300,1.00000000)
 (2.7600,1.00000000) (2.7900,1.00000000) (2.8200,1.00000000) (2.8500,1.00000000) (2.8800,1.00000000) (2.9100,1.00000000) (2.9400,1.00000000) (2.9700,1.00000000)
 (3.0000,1.00000000)
};
\draw[AnalysisGray,thick,densely dashed] plot coordinates {
 (-3.0000,0.00001105) (-2.9700,0.00001333) (-2.9400,0.00001607) (-2.9100,0.00001933) (-2.8800,0.00002321) (-2.8500,0.00002783) (-2.8200,0.00003330) (-2.7900,0.00003979)
 (-2.7600,0.00004746) (-2.7300,0.00005651) (-2.7000,0.00006717) (-2.6700,0.00007970) (-2.6400,0.00009441) (-2.6100,0.00011164) (-2.5800,0.00013180) (-2.5500,0.00015533)
 (-2.5200,0.00018275) (-2.4900,0.00021464) (-2.4600,0.00025168) (-2.4300,0.00029460) (-2.4000,0.00034426) (-2.3700,0.00040161) (-2.3400,0.00046772) (-2.3100,0.00054379)
 (-2.2800,0.00063117) (-2.2500,0.00073136) (-2.2200,0.00084603) (-2.1900,0.00097703) (-2.1600,0.00112642) (-2.1300,0.00129649) (-2.1000,0.00148973) (-2.0700,0.00170892)
 (-2.0400,0.00195710) (-2.0100,0.00223758) (-1.9800,0.00255400) (-1.9500,0.00291033) (-1.9200,0.00331089) (-1.8900,0.00376034) (-1.8600,0.00426376) (-1.8300,0.00482660)
 (-1.8000,0.00545475) (-1.7700,0.00615453) (-1.7400,0.00693270) (-1.7100,0.00779650) (-1.6800,0.00875361) (-1.6500,0.00981221) (-1.6200,0.01098096) (-1.5900,0.01226899)
 (-1.5600,0.01368594) (-1.5300,0.01524190) (-1.5000,0.01694743) (-1.4700,0.01881355) (-1.4400,0.02085172) (-1.4100,0.02307378) (-1.3800,0.02549198) (-1.3500,0.02811890)
 (-1.3200,0.03096742) (-1.2900,0.03405068) (-1.2600,0.03738203) (-1.2300,0.04097495) (-1.2000,0.04484301) (-1.1700,0.04899980) (-1.1400,0.05345884) (-1.1100,0.05823349)
 (-1.0800,0.06333692) (-1.0500,0.06878195) (-1.0200,0.07458099) (-0.9900,0.08074597) (-0.9600,0.08728818) (-0.9300,0.09421822) (-0.9000,0.10154589) (-0.8700,0.10928008)
 (-0.8400,0.11742864) (-0.8100,0.12599836) (-0.7800,0.13499478) (-0.7500,0.14442218) (-0.7200,0.15428344) (-0.6900,0.16457997) (-0.6600,0.17531166) (-0.6300,0.18647678)
 (-0.6000,0.19807195) (-0.5700,0.21009210) (-0.5400,0.22253037) (-0.5100,0.23537819) (-0.4800,0.24862516) (-0.4500,0.26225914) (-0.4200,0.27626619) (-0.3900,0.29063065)
 (-0.3600,0.30533515) (-0.3300,0.32036067) (-0.3000,0.33568662) (-0.2700,0.35129089) (-0.2400,0.36714997) (-0.2100,0.38323904) (-0.1800,0.39953208) (-0.1500,0.41600201)
 (-0.1200,0.43262082) (-0.0900,0.44935970) (-0.0600,0.46618920) (-0.0300,0.48307939) (0.0000,0.50000000) (0.0300,0.51692061) (0.0600,0.53381080) (0.0900,0.55064030)
 (0.1200,0.56737918) (0.1500,0.58399799) (0.1800,0.60046792) (0.2100,0.61676096) (0.2400,0.63285003) (0.2700,0.64870911) (0.3000,0.66431338) (0.3300,0.67963933)
 (0.3600,0.69466485) (0.3900,0.70936935) (0.4200,0.72373381) (0.4500,0.73774086) (0.4800,0.75137484) (0.5100,0.76462181) (0.5400,0.77746963) (0.5700,0.78990790)
 (0.6000,0.80192805) (0.6300,0.81352322) (0.6600,0.82468834) (0.6900,0.83542003) (0.7200,0.84571656) (0.7500,0.85557782) (0.7800,0.86500522) (0.8100,0.87400164)
 (0.8400,0.88257136) (0.8700,0.89071992) (0.9000,0.89845411) (0.9300,0.90578178) (0.9600,0.91271182) (0.9900,0.91925403) (1.0200,0.92541901) (1.0500,0.93121805)
 (1.0800,0.93666308) (1.1100,0.94176651) (1.1400,0.94654116) (1.1700,0.95100020) (1.2000,0.95515699) (1.2300,0.95902505) (1.2600,0.96261797) (1.2900,0.96594932)
 (1.3200,0.96903258) (1.3500,0.97188110) (1.3800,0.97450802) (1.4100,0.97692622) (1.4400,0.97914828) (1.4700,0.98118645) (1.5000,0.98305257) (1.5300,0.98475810)
 (1.5600,0.98631406) (1.5900,0.98773101) (1.6200,0.98901904) (1.6500,0.99018779) (1.6800,0.99124639) (1.7100,0.99220350) (1.7400,0.99306730) (1.7700,0.99384547)
 (1.8000,0.99454525) (1.8300,0.99517340) (1.8600,0.99573624) (1.8900,0.99623966) (1.9200,0.99668911) (1.9500,0.99708967) (1.9800,0.99744600) (2.0100,0.99776242)
 (2.0400,0.99804290) (2.0700,0.99829108) (2.1000,0.99851027) (2.1300,0.99870351) (2.1600,0.99887358) (2.1900,0.99902297) (2.2200,0.99915397) (2.2500,0.99926864)
 (2.2800,0.99936883) (2.3100,0.99945621) (2.3400,0.99953228) (2.3700,0.99959839) (2.4000,0.99965574) (2.4300,0.99970540) (2.4600,0.99974832) (2.4900,0.99978536)
 (2.5200,0.99981725) (2.5500,0.99984467) (2.5800,0.99986820) (2.6100,0.99988836) (2.6400,0.99990559) (2.6700,0.99992030) (2.7000,0.99993283) (2.7300,0.99994349)
 (2.7600,0.99995254) (2.7900,0.99996021) (2.8200,0.99996670) (2.8500,0.99997217) (2.8800,0.99997679) (2.9100,0.99998067) (2.9400,0.99998393) (2.9700,0.99998667)
 (3.0000,0.99998895)
};
\draw[black!55,thick,dash dot] plot coordinates {
 (-3.0000,0.01694743) (-2.9700,0.01785971) (-2.9400,0.01881355) (-2.9100,0.01981040) (-2.8800,0.02085172) (-2.8500,0.02193900) (-2.8200,0.02307378) (-2.7900,0.02425759)
 (-2.7600,0.02549198) (-2.7300,0.02677855) (-2.7000,0.02811890) (-2.6700,0.02951465) (-2.6400,0.03096742) (-2.6100,0.03247888) (-2.5800,0.03405068) (-2.5500,0.03568451)
 (-2.5200,0.03738203) (-2.4900,0.03914494) (-2.4600,0.04097495) (-2.4300,0.04287374) (-2.4000,0.04484301) (-2.3700,0.04688447) (-2.3400,0.04899980) (-2.3100,0.05119070)
 (-2.2800,0.05345884) (-2.2500,0.05580588) (-2.2200,0.05823349) (-2.1900,0.06074330) (-2.1600,0.06333692) (-2.1300,0.06601595) (-2.1000,0.06878195) (-2.0700,0.07163646)
 (-2.0400,0.07458099) (-2.0100,0.07761701) (-1.9800,0.08074597) (-1.9500,0.08396924) (-1.9200,0.08728818) (-1.8900,0.09070409) (-1.8600,0.09421822) (-1.8300,0.09783178)
 (-1.8000,0.10154589) (-1.7700,0.10536165) (-1.7400,0.10928008) (-1.7100,0.11330211) (-1.6800,0.11742864) (-1.6500,0.12166048) (-1.6200,0.12599836) (-1.5900,0.13044294)
 (-1.5600,0.13499478) (-1.5300,0.13965440) (-1.5000,0.14442218) (-1.4700,0.14929845) (-1.4400,0.15428344) (-1.4100,0.15937726) (-1.3800,0.16457997) (-1.3500,0.16989149)
 (-1.3200,0.17531166) (-1.2900,0.18084021) (-1.2600,0.18647678) (-1.2300,0.19222089) (-1.2000,0.19807195) (-1.1700,0.20402929) (-1.1400,0.21009210) (-1.1100,0.21625946)
 (-1.0800,0.22253037) (-1.0500,0.22890370) (-1.0200,0.23537819) (-0.9900,0.24195250) (-0.9600,0.24862516) (-0.9300,0.25539461) (-0.9000,0.26225914) (-0.8700,0.26921697)
 (-0.8400,0.27626619) (-0.8100,0.28340479) (-0.7800,0.29063065) (-0.7500,0.29794155) (-0.7200,0.30533515) (-0.6900,0.31280903) (-0.6600,0.32036067) (-0.6300,0.32798744)
 (-0.6000,0.33568662) (-0.5700,0.34345540) (-0.5400,0.35129089) (-0.5100,0.35919010) (-0.4800,0.36714997) (-0.4500,0.37516736) (-0.4200,0.38323904) (-0.3900,0.39136173)
 (-0.3600,0.39953208) (-0.3300,0.40774667) (-0.3000,0.41600201) (-0.2700,0.42429459) (-0.2400,0.43262082) (-0.2100,0.44097708) (-0.1800,0.44935970) (-0.1500,0.45776499)
 (-0.1200,0.46618920) (-0.0900,0.47462860) (-0.0600,0.48307939) (-0.0300,0.49153779) (0.0000,0.50000000) (0.0300,0.50846221) (0.0600,0.51692061) (0.0900,0.52537140)
 (0.1200,0.53381080) (0.1500,0.54223501) (0.1800,0.55064030) (0.2100,0.55902292) (0.2400,0.56737918) (0.2700,0.57570541) (0.3000,0.58399799) (0.3300,0.59225333)
 (0.3600,0.60046792) (0.3900,0.60863827) (0.4200,0.61676096) (0.4500,0.62483264) (0.4800,0.63285003) (0.5100,0.64080990) (0.5400,0.64870911) (0.5700,0.65654460)
 (0.6000,0.66431338) (0.6300,0.67201256) (0.6600,0.67963933) (0.6900,0.68719097) (0.7200,0.69466485) (0.7500,0.70205845) (0.7800,0.70936935) (0.8100,0.71659521)
 (0.8400,0.72373381) (0.8700,0.73078303) (0.9000,0.73774086) (0.9300,0.74460539) (0.9600,0.75137484) (0.9900,0.75804750) (1.0200,0.76462181) (1.0500,0.77109630)
 (1.0800,0.77746963) (1.1100,0.78374054) (1.1400,0.78990790) (1.1700,0.79597071) (1.2000,0.80192805) (1.2300,0.80777911) (1.2600,0.81352322) (1.2900,0.81915979)
 (1.3200,0.82468834) (1.3500,0.83010851) (1.3800,0.83542003) (1.4100,0.84062274) (1.4400,0.84571656) (1.4700,0.85070155) (1.5000,0.85557782) (1.5300,0.86034560)
 (1.5600,0.86500522) (1.5900,0.86955706) (1.6200,0.87400164) (1.6500,0.87833952) (1.6800,0.88257136) (1.7100,0.88669789) (1.7400,0.89071992) (1.7700,0.89463835)
 (1.8000,0.89845411) (1.8300,0.90216822) (1.8600,0.90578178) (1.8900,0.90929591) (1.9200,0.91271182) (1.9500,0.91603076) (1.9800,0.91925403) (2.0100,0.92238299)
 (2.0400,0.92541901) (2.0700,0.92836354) (2.1000,0.93121805) (2.1300,0.93398405) (2.1600,0.93666308) (2.1900,0.93925670) (2.2200,0.94176651) (2.2500,0.94419412)
 (2.2800,0.94654116) (2.3100,0.94880930) (2.3400,0.95100020) (2.3700,0.95311553) (2.4000,0.95515699) (2.4300,0.95712626) (2.4600,0.95902505) (2.4900,0.96085506)
 (2.5200,0.96261797) (2.5500,0.96431549) (2.5800,0.96594932) (2.6100,0.96752112) (2.6400,0.96903258) (2.6700,0.97048535) (2.7000,0.97188110) (2.7300,0.97322145)
 (2.7600,0.97450802) (2.7900,0.97574241) (2.8200,0.97692622) (2.8500,0.97806100) (2.8800,0.97914828) (2.9100,0.98018960) (2.9400,0.98118645) (2.9700,0.98214029)
 (3.0000,0.98305257)
};
\draw[AnalysisGray,thin,densely dashed] (-0.95387255,-.24)--(-0.95387255,-.16)
 (-0.95387255,.02)--(-0.95387255,.25);
\draw[AnalysisGray,thin,densely dashed] (0.95387255,-.24)--(0.95387255,-.16)
 (0.95387255,.02)--(0.95387255,.75);
\fill[black!55] (-0.95387255,.25) circle(1.2pt);
\fill[black!55] (0.95387255,.75) circle(1.2pt);
\draw[<->,AnalysisGray] (-0.95387255,-.24)--(0.95387255,-.24);
\node[align=center,font=\small] at (0,-.39){固定温度比例间的过渡宽度\(\asymp\sqrt{\kappa t}\)};
\draw[black,thick,densely dotted] (3.9,1.02)--(4.5,1.02);
\node[anchor=west,font=\small] at (4.6,1.02){\(t=0\)};
\draw[AnalysisBlue,very thick] (3.9,.82)--(4.5,.82);
\node[anchor=west,font=\small] at (4.6,.82){\(t=1/16\)};
\draw[AnalysisGray,thick,densely dashed] (3.9,.62)--(4.5,.62);
\node[anchor=west,font=\small] at (4.6,.62){\(t=1/4\)};
\draw[black!55,thick,dash dot] (3.9,.42)--(4.5,.42);
\node[anchor=west,font=\small] at (4.6,.42){\(t=1\)};
\end{tikzpicture}
\caption[阶跃初温的光滑化与扩散尺度]{阶跃初温的光滑化与扩散尺度。图取\(\kappa=1\)，点线为\(f=\mathbf1_{(0,\infty)}\)，其余三条曲线按正文的 Gaussian 积分计算；下方标出\(t=1\)时温度为\(1/4\)、\(3/4\)的两点间距。任意固定温度比例间的过渡宽度与\(\sqrt{\kappa t}\)成正比。每个正时间剖面均光滑，且对每个有限\(x\)都有\(0<u(t,x)<1\)：原先零温的任意远处已出现正温度，但远离界面的影响可以极弱，图上线条与水平线近乎重合不表示尾部为零。}
\label{b2:fig:12-heat-step}
\end{figure}


\subsection{最大值原理}
\label{b2:subsec:120302}

最大值原理从微分方程本身控制解，不依赖可显式计算的核。本小节讨论实值函数。若 \(\Omega\subset\mathbb R^d\) 是有界开集、\(T>0\)，记 \(Q_T=(0,T]\times\Omega\)，其\term[paowu-bianjie]{抛物边界}[parabolic boundary]为
\[
 \partial_pQ_T=(\{0\}\times\overline\Omega)
 \cup([0,T]\times\partial\Omega).
\]
上端面 \(t=T\) 的内部点不属于这条边界：它们仍受到先前时间的数据控制。

\Needspace{6\baselineskip}
\begin{theorem}[抛物最大值原理]\label{b2:thm:12-maximum}
设整数 \(d\geq1\)、\(\kappa,T>0\)，\(\Omega\subset\mathbb R^d\) 为非空有界开集，\(\partial_pQ_T=(\{0\}\times\overline\Omega)\cup([0,T]\times\partial\Omega)\)。设实值 \(u\in C([0,T]\times\overline\Omega)\)，在正时间内有一阶时间导数和二阶空间导数，且 \(u_t-\kappa\Delta u\leq0\)。则
\[
 \max_{[0,T]\times\overline\Omega}u
 \leq\max_{\partial_pQ_T}u.
\]
导数在 \(t=T\) 取左侧时间导数即可。
\end{theorem}
\begin{proof}
固定 \(\varepsilon>0\)，令 \(v=u-\varepsilon t\)，则 \(v_t-\kappa\Delta v<0\)。若 \(v\) 的最大值严格超过抛物边界的最大值，紧性使它在某个正时间、空间内部的点取得最大值。该点有 \(\nabla_xv=0\)、\(\Delta v\leq0\)，且相对于过去时间有 \(v_t\geq0\)，故 \(v_t-\kappa\Delta v\geq0\)，矛盾。于是 \(u(t,x)\leq\max_{\partial_pQ_T}u+\varepsilon T\)，令 \(\varepsilon\downarrow0\)。
\end{proof}

这称为\term[paowu-zuidazhi-yuanli]{抛物最大值原理}[parabolic maximum principle]。对 \(-u\) 使用它得到最小值控制，对两个解之差使用它得到比较原理。特别是同样的 Dirichlet 边界条件下，初值误差的一致范数不会增长。有非齐次源项时，若 \(|F|\leq M\)，比较函数可加上 \(Mt\)，从而得到初边值最大误差再加 \(MT\) 的估计。

\ProseParagraphBreak

无界空间不能直接援引紧性取得最大值，也不能忽略远处边界。对每个有限时间带上一致有界的经典次解，可以用障碍函数
\[
 \Phi(t,x)=1+|x|^2+2d\kappa t,
 \quad\Phi_t-\kappa\Delta\Phi=0
\]
补上这一步。若初值 \(u(0,x)\leq0\)，在半径 \(R\) 的球柱中比较 \(u\) 与 \(\varepsilon\Phi\)。有界性使足够大的侧边界满足 \(u\leq\varepsilon\Phi\)，初始面也满足；有界区域最大值原理于是给出任意固定点的同一不等式。先令 \(R\) 足够大，再令 \(\varepsilon\downarrow0\)，得到 \(u\leq0\)。

\subsection{唯一性}
\label{b2:subsec:120303}

上面的障碍证明给出了全空间有界经典解的唯一性：在每个 \([0,T]\times\mathbb R^d\) 上有共同有限界、连续取得相同初值的两个解必相等。该结论不覆盖任意增长的经典解。空间无穷远的限制参与证明，不能从有界区域的唯一性直接删去。

\ProseParagraphBreak

对平方可积初值，另一种自然解类由能量决定。以下版本只要求分布方程与时间上的 \(L^2\) 连续性，因此也覆盖此前构造的粗初值解。

\Needspace{6\baselineskip}
\begin{theorem}[平方可积解的唯一性]\label{b2:thm:12-l2-uniqueness}
设整数 \(d\geq1\)、\(\kappa>0\)，\(H_s(x)=(4\pi s)^{-d/2}e^{-|x|^2/(4s)}\)（\(s>0\)、\(x\in\mathbb R^d\)）。设 \(T>0\)、\(f\in L^2(\mathbb R^d;\mathbb C)\)。若 \(u\in C([0,T];L^2(\mathbb R^d))\) 在分布意义下满足 \(u_t=\kappa\Delta u\)，且 \(u(0)=f\)，则 \(u(t)=H_{\kappa t}*f\)。此外，核解满足
\begin{equation}\label{b2:eq:12-whole-energy}
 \|u(T)\|_2^2+2\kappa\int_0^T\|\nabla u(t)\|_2^2\,\mathrm{d}t
 =\|f\|_2^2.
\end{equation}
\end{theorem}
\begin{proof}
空间 Fourier 变换在 \(L^2\) 上酉。先说明可以对方程使用 Schwartz 空间检验函数：将它乘以趋于一的空间紧支集截断，函数及二阶导数都在 \(L^2\) 中收敛；\(u\) 在紧时间区间上范数有界，所以分布恒等式传到极限。因此空间变换后有 \(\partial_t\widehat u+a\widehat u=0\)，其中 \(a(\xi)=4\pi^2\kappa|\xi|^2\)。

取任意紧支集光滑频率函数 \(\chi\)。在其支集上，\(e^{ta(\xi)}\) 及所需导数有界。乘以这个因子得到
\[
 \partial_t\left(e^{ta}\chi\widehat u(t)\right)=0.
\]
对每个光滑紧支集的频率测试函数配对，所得标量分布导数为零，故是常数；它又随时间连续，常数由初值确定。因此 \(\chi\widehat u(t)=e^{-ta}\chi\widehat f\)。取覆盖所有频率的截断列，得到 \(\widehat u(t)=e^{-ta}\widehat f\) 于 \(L^2\)，从而唯一。

最后对每个频率有
\[
 e^{-2Ta}|\widehat f|^2
 +2\int_0^Tae^{-2ta}|\widehat f|^2\,\mathrm{d}t
 =|\widehat f|^2.
\]
非负积分可换序；Plancherel 与平方可积弱导数的频率刻画把第二项变为 \(2\kappa\int\|\nabla u\|_2^2\)，得到能量恒等式。
\end{proof}

有界经典类与 \(C_tL^2_x\) 类各有独立的远处控制：前者允许非零常数初温，后者容许无界但平方可积的初温。二者有交集，却不相互包含。对一个给定初值问题，先明确解类，再谈唯一性。

\Needspace{6\baselineskip}
\begin{proposition}[粗外力与热能量]\label{b2:prop:12-forced-energy}
设整数 \(d\geq1\)、\(\kappa>0\)，\(H_s(x)=(4\pi s)^{-d/2}e^{-|x|^2/(4s)}\)（\(s>0\)、\(x\in\mathbb R^d\)）。设 \(T>0\)、\(f\in L^2(\mathbb R^d;\mathbb C)\)、\(F\in L^1((0,T);L^2(\mathbb R^d))\)。式\eqref{b2:eq:12-duhamel}的积分以 \(L^2\) 值积分解释，给出唯一的 \(C([0,T];L^2)\) 分布解，满足 \(u_t-\kappa\Delta u=F\) 及初值 \(f\)。它满足
\[
 \|u(t)\|_2\leq\|f\|_2+\int_0^t\|F(s)\|_2\,\mathrm{d}s,
 \quad \nabla u\in L^2((0,T)\times\mathbb R^d),
\]
以及能量恒等式
\[
 \|u(t)\|_2^2+2\kappa\int_0^t\|\nabla u(s)\|_2^2\,\mathrm{d}s
 =\|f\|_2^2+2\operatorname{Re}\int_0^t\langle F(s),u(s)\rangle_{L^2}\,\mathrm{d}s.
\]
\end{proposition}
\begin{proof}
热卷积在 \(L^2\) 中为收缩，并在零时强趋于恒等。先对简单外力定义积分，再以积分范数逼近，得到公式及第一项估计。时间连续性可先对有限个空间取值的简单函数检验：各段积分由核的强连续性和控制收敛得到连续，端点处新增积分的范数趋于零。一般外力的误差在整个 \([0,T]\) 上不超过其 \(L^1_tL^2_x\) 范数，故同样连续，且积分项在零时趋于零。

选光滑紧支集 \(f_j\to f\) 于 \(L^2\)，并选光滑时空外力 \(F_j\to F\) 于 \(L^1_tL^2_x\)，使每个固定 \(j\) 的 \(F_j\) 在整个时间带上有紧空间支集，紧集允许随 \(j\) 扩大。后者可先由简单函数逼近，再分别磨光有限个空间取值及时间示性函数。对应光滑解 \(u_j\) 在 \(C_tL^2_x\) 中趋于 \(u\)，方程的测试等式随之传到极限。对这些光滑解乘以 \(\overline{u_j}\)、取实部并作空间积分，得到所称能量恒等式。全空间分部积分可先乘标准截断 \(\chi_R^2\)：额外项的模至多为 \(CR^{-1}\|u_j\|_2\cdot\|\nabla u_j\|_2\)。对每个固定 \(j\)，初值与外力的光滑紧支集逼近使 \(\|u_j(t)\|_2\)、\(\|\nabla u_j(t)\|_2\) 在该有限时间带上有界，这也可将一次空间导数移到平滑初值和外力上后用 Young 不等式得到。因此截断通量的时间积分趋于零，再令 \(R\to\infty\) 即得上述恒等式；这里不要求对所有 \(j\) 共用一个空间支集。

还需证明梯度收敛，不能仅由零阶收敛推断。记两近似解之差为 \(v\)，初值差为 \(h\)，外力差为 \(G\)。差解的能量恒等式给出
\[
 2\kappa\int_0^T\|\nabla v(s)\|_2^2\,\mathrm{d}s
 \leq\|h\|_2^2+2\sup_{s\leq T}\|v(s)\|_2
                    \int_0^T\|G(s)\|_2\,\mathrm{d}s.
\]
第一项估计又控制 \(\sup\|v\|_2\)，所以梯度在时空 \(L^2\) 中为 Cauchy 列。对空间测试函数分部积分取极限，将其极限识别为 \(\nabla u\)。于是能量的耗散积分收敛；外力功由外力的积分范数收敛及 \(u_j\) 的一致 \(L^2\) 收敛传到极限。所得等式对每个 \(t\in[0,T]\) 成立。两个同数据解之差满足齐次方程及零初值，\cref{b2:thm:12-l2-uniqueness}给出唯一性。
\end{proof}

\subsection{保正性}
\label{b2:subsec:120304}

设 \(\kappa>0\)、\(1\le p\le\infty\)、\(f\in L^p(\mathbb R^d;\mathbb R)\)。热核严格为正。因此若 \(f\geq0\) 几乎处处，则 \(H_{\kappa t}*f\geq0\)；若 \(f\) 不是零等价类，某个正测度集合上有 \(f>0\)，对每个 \(t>0,x\) 的积分都严格为正。这一性质称为\term[baozheng-xing]{保正性}[positivity preservation]。即使初温支集很小，任何正时间全空间每一点都受到影响；热传播没有波方程的有限传播速度。

\ProseParagraphBreak

在有限区间 \((0,\ell)\)、\(\ell,\kappa>0\) 上，纯热方程 \(u_t=\kappa u_{xx}\) 的非负实初温 \(f\in L^2(0,\ell;\mathbb R)\) 也产生非负解。对零 Dirichlet 条件，先用非负光滑紧支集函数逼近初值，最大值原理给出各近似解非负；正时间核的 \(L^2\to L^\infty\) 界让它们一致趋于目标解，故极限非负。零 Neumann 的镜像核 \(K_N\) 本身为正。Robin 非负系数时，可以直接使用负部能量，沿用\secref{b2:sec:1201}的通量约定，边界写为 \(-\kappa u_x(0)+h_0u(0)=\kappa u_x(\ell)+h_\ell u(\ell)=0\)，\(h_0,h_\ell\geq0\)。正时间有 \(u(t)\in D(A)\)、\(u_t(t)\in H\)，因此不需预设初温有经典导数。

\ProseParagraphBreak

令 \(u_-=\max\{-u,0\}\)。绝对连续函数与这个 Lipschitz 函数复合后仍绝对连续，且 \((u_-)'=-u'\cdot\mathbf1_{\{u<0\}}\) 几乎处处；在零水平集上 \(u'=0\) 几乎处处，这可在可微且为该水平集密度点的位置沿零点差商取得。故 \(u_-\in V\)。对标量函数 \(J(z)=z_-^2/2\)，导数为 \(-z_-\)，且导数为 Lipschitz 函数，因此 Taylor 余项被 \(|h|^2/2\) 控制。积分这个余项，再使用 \(H\) 中的时间可微性，得到 \(\dfrac{\mathrm{d}}{\mathrm{d}t}\int J(u)=-\langle u_t,u_-\rangle\)。将方程与 \(u_-\) 配对，并使用 \(\mathfrak a(u,u_-)=-\kappa\|(u_-)'\|_2^2-\sum_s h_s|u_-(s)|^2\)，便有
\[
 \dfrac12\dfrac{\mathrm{d}}{\mathrm{d}t}\|u_-\|_2^2
 =-\kappa\|(u_-)'\|_2^2
 -\sum_{s\in R}h_s|u_-(s)|^2\leq0.
\]
在 \([\varepsilon,t]\) 上积分，得到 \(\|u_-(t)\|_2\leq\|u_-(\varepsilon)\|_2\)。负部映射在 \(L^2\) 中是 Lipschitz 的，而 \(u(\varepsilon)\to f\) 于 \(L^2\)，令 \(\varepsilon\downarrow0\) 得 \(\|u_-(t)\|_2\leq\|f_-\|_2\)。零初始负部于是始终为零。论证只使用一次空间弱导数，不要求负部在零点具有经典二阶导数。

\ProseParagraphBreak

全空间的纯热核解还保持总质量。若 \(f\in L^1\)，由绝对可积换序有 \(\int u(t)=\int f\)。对非负 \(f\)，当 \(0<s<t\) 时，半群关系及热核非负、总质量为一给出 \(\|u(t)\|_\infty\leq\|u(s)\|_\infty\)。降低的是全空间的峰值，不能据此说每一点的温度都随时间下降：初温支集外的零温点在正时间会升温。热量在空间中重新分布，而总量保持不变。若初值还有有限二阶绝对矩，Gauss 核的零一阶矩与每方向二阶矩 \(2\kappa t\) 给出
\[
 \int|x|^2u(t,x)\,\mathrm{d}x
 =\int|y|^2f(y)\,\mathrm{d}y+2d\kappa t\int f(y)\,\mathrm{d}y.
\]
当 \(f\geq0\) 时，这是总热量向半径 \(\sqrt t\) 的尺度扩散的定量表达。

% !TeX root = ../../Book2.tex
\section{热半群}
\label{b2:sec:1204}

\subsection{半群性质}
\label{b2:subsec:120401}

将时刻 \(t\) 的解视为对初值的作用，记
\[
 S(t)f\coloneqq H_{\kappa t}*f\quad(t>0),
 \quad S(0)f=f.
\]
这一族算子称为\term[re-banqun]{热半群}[heat semigroup]。先演化 \(s\) 再演化 \(t\)，应与直接演化 \(s+t\) 相同；核公式使这一预期成为可以验证的等式。
\symidx{\(S(t)\)：全空间热半群}

\begin{proposition}\label{b2:prop:12-semigroup}
对 \(s,t\geq0\)，有 \(S(t)S(s)=S(t+s)\)。当 \(s,t>0\) 时，等价的核关系为 \(H_{\kappa t}*H_{\kappa s}=H_{\kappa(t+s)}\)。这些等式在全部 \(L^p\)、\(1\leq p\leq\infty\)，以及有界一致连续函数空间上成立。
\end{proposition}
\begin{proof}
两个热核均可积，卷积定理给出变换
\[
 e^{-4\pi^2\kappa t|\xi|^2}e^{-4\pi^2\kappa s|\xi|^2}
 =e^{-4\pi^2\kappa(t+s)|\xi|^2}.
\]
Fourier 唯一性使核相等。对一般 \(L^p\) 初值，卷积结合律可先对有界紧支集函数由 Fubini 证明，再由 Young 不等式和稠密性推广；\(p=\infty\) 时，双重积分的绝对值直接由 \(\|f\|_\infty\) 乘两个核的质量控制，Fubini 仍适用。零参数时是恒等算子。
\end{proof}

区间上的对应算子记为 \(S_I(t)\)，由式\eqref{b2:eq:12-interval-expansion}按每个特征系数相乘，也有同一半群关系。全空间借助平移不变的核，区间借助边界适配的特征基；两者组织时间的方式相同。

\subsection{压缩性}
\label{b2:subsec:120402}

\begin{proposition}[范数收缩与强连续性]\label{b2:prop:12-contraction}
全空间热半群满足 \(\|S(t)f\|_p\leq\|f\|_p\)，\(1\leq p\leq\infty\)。对 \(p<\infty\)，映射 \(t\mapsto S(t)f\) 在 \([0,\infty)\) 上按 \(L^p\) 范数连续。这个连续性也在 \(C_0(\mathbb R^d)\) 与有界一致连续函数空间上成立，但不在整个 \(L^\infty\) 上成立。
\end{proposition}
\begin{proof}
热核质量为一，Young 不等式给出范数界。零时连续性已由逼近恒等证明。若 \(h>0\)，半群关系给出
\[
 \|S(t+h)f-S(t)f\|_p\leq\|S(h)f-f\|_p.
\]
对左侧时间差，用 \(S(t)f-S(t-h)f=S(t-h)(S(h)f-f)\) 得同样控制。于是零时连续性延伸到全部时间。

对 \(C_0\) 还须检查不变性。连续且在无穷远消失的函数可由连续紧支集函数一致逼近。后者与热核卷积在无穷远消失：初值支集有界，核在 \(|x|\to\infty\) 时于该支集上一致趋零。范数界让这项性质传到一般 \(C_0\) 函数；连续性和一致连续性由核的平移连续性得到。有界一致连续空间的不变性则直接由 \(|S(t)f(x+h)-S(t)f(x)|\leq\sup_y|f(y+h)-f(y)|\) 得出。

最后取 \(f=\mathbf1_{\{x_1>0\}}\)。一维积分公式表明对每个 \(t>0\)，\(\|S(t)f-f\|_\infty=1/2\)：在界面两侧趋近界面，误差趋于 \(1/2\)，而远离界面时严格更小。因此全体 \(L^\infty\) 的强连续性失败。
\end{proof}

范数不增不等于每个非零函数都按一个共同的指数衰减。在 \(L^2(\mathbb R^d)\) 上，乘子 \(e^{-4\pi^2\kappa t|\xi|^2}\) 的本性上确界为一，因此 \(\|S(t)\|_{2\to2}=1\)。任意低频附近的数据衰减都可以很慢。另一方面，对每个固定 \(f\in L^2\)，这个乘子在几乎每个频率处趋零，控制收敛给出 \(\|S(t)f\|_2\to0\)。逐个初值衰减与统一算子范数衰减不能混为一谈。

\ProseParagraphBreak

Dirichlet 与 Neumann 区间热算子也在全部 \(L^p\) 上收缩。Neumann 核非负且对两变量积分都为一；Dirichlet 核非负且两类积分均不超过一。后一个断言可由初值一的解与常数一比较得到，再用核对称性换变量。对 \(1<p<\infty\)，补上一份零函数的质量后应用 Jensen，不等式为 \(|S_I(t)f|^p\leq S_I(t)|f|^p\)，再积分即得；两个端点直接估计。Dirichlet 核的非负性由上一节正性保持与连续核得到，若某个核值为负，以支集集中于该点附近的非负初值测试就会矛盾。

\subsection{生成元观点}
\label{b2:subsec:120403}

半群给出了有限时间的变化，生成元记录零时的一阶变化。设 \((T(t))_{t\geq0}\) 是 Banach 空间 \(X\) 上满足 \(T(0)=I\)、\(T(t+s)=T(t)T(s)\) 的有界线性算子族，并且 \(T(t)x\to x\) 当 \(t\downarrow0\) 对每个 \(x\in X\) 成立。这样的族称为\term[qiang-lianxu-banqun]{强连续半群}[strongly continuous semigroup]。定义
\begin{equation}\label{b2:eq:12-generator-def}
 D(L)\coloneqq\left\{\,f\in X\;\middle|\;
 \lim_{t\downarrow0}\dfrac{T(t)f-f}{t}\text{ 在 }X\text{ 中存在}\,\right\},
 \quad Lf\coloneqq\lim_{t\downarrow0}\dfrac{T(t)f-f}{t},
\end{equation}
称 \(L\) 为半群的\term[shengchengyuan]{生成元}[generator]。这个极限必须在所选空间的范数中成立，不能用逐点导数替代定义域条件。

\begin{theorem}[热半群的平方可积生成元]\label{b2:thm:12-generator-l2}
在 \(X=L^2(\mathbb R^d)\) 上，热半群的生成元为
\[
 L=\kappa\Delta,
 \quad D(L)=H^2(\mathbb R^d).
\]
对 \(\lambda>0\)，其逆算子有
\[
 (\lambda I-L)^{-1}f=\int_0^\infty e^{-\lambda t}S(t)f\,\mathrm{d}t,
 \quad\|(\lambda I-L)^{-1}\|_{2\to2}\leq\lambda^{-1}.
\]
\end{theorem}
\begin{proof}
仍令 \(a(\xi)=4\pi^2\kappa|\xi|^2\)。若 \(a\widehat f\in L^2\)，则
\[
 \left|\dfrac{e^{-ta}-1}{t}\right|\leq a,
 \quad\dfrac{e^{-ta}-1}{t}\widehat f\longrightarrow-a\widehat f
 \quad\text{于 }L^2.
\]
控制收敛与 Plancherel 说明 \(f\in D(L)\)，且 \(Lf=\kappa\Delta f\)。反过来，若差商在 \(L^2\) 中收敛，变换也收敛；取一个几乎处处收敛的子列，与乘子逐点极限比较，得到极限必为 \(-a\widehat f\)，故 \(a\widehat f\in L^2\)。由\cref{b2:thm:04-weak-derivative}，条件 \(f,|\xi|^2\widehat f\in L^2\) 等价于 \(f\in H^2\)：各二次单项式均受 \(|\xi|^2\) 控制，一次单项式则由 \(1+|\xi|^2\) 控制。

最后，指数权重使算子积分在 \(L^2\) 范数中收敛。变换后可逐频率积分，乘子为 \(\int_0^\infty e^{-t(\lambda+a)}\,\mathrm{d}t=(\lambda+a)^{-1}\)。其模不超过 \(\lambda^{-1}\)，而 \(a/(\lambda+a)\leq1\)，故积分的像属于生成元定义域。乘上 \(\lambda+a\) 即验证双侧逆。
\end{proof}

区间的生成元为 \(-A\)，定义域仍是\cref{b2:def:10-sl-operator}的边界定义域；系数差商与同一控制论证证明这一点。负号由 \(u_t+Au=0\) 决定。相同的微分表达式 \(\kappa\partial_x^2\)，配上不同边界定义域，就产生不同的半群和不同的长期平衡。

\subsection{\texorpdfstring{\(C_0\)}{C0} 半群的预告}
\label{b2:subsec:120404}

强连续半群也称 \(C_0\) 半群；这里的下标表示零时强连续，与函数空间 \(C_0(\mathbb R^d)\) 表示无穷远消失不是同一个定义。若还满足 \(\|T(t)\|\leq1\)，称为\term[yasuo-banqun]{压缩半群}[contraction semigroup]；此处允许范数等于一，不要求压缩映射不动点定理中的严格压缩常数。

\ProseParagraphBreak

在热流中，半群先由显式核构造，再识别生成元。更一般的扩散模型常没有可写出的核，问题顺序会反过来：先给闭算子及其定义域，再由解算子的估计构造半群。\appref{b2:app:D}将介绍这条抽象路线。当前结果只需读者掌握三点：时间连续性依赖所选空间，生成元必须连同定义域一起给出，指数演化的记号并不自动允许反向时间。

\ProseParagraphBreak

反向时间会放大高频。在全空间 \(L^2\) 中，若想由时刻 \(T>0\) 的值恢复初值，频率公式要求乘以 \(e^{4\pi^2\kappa T|\xi|^2}\)。该乘子无界，故即使正时间演化是单射，逆问题也不具有同一范数下的连续依赖。温度演化的稳定性与反演初温的稳定性是两个不同的问题。

% 本章小结（不计入编号）
\begin{chaptersummary}
区间热流由边界适配的模态展开，全空间热流由 Gauss 卷积表示；两种构造都能从粗初值出发，通过范数极限取得初值。正时间的指数频率衰减提供光滑化，但零模态是否保留、总热量是否守恒、是否有统一指数衰减，都要由所在区域和边界条件判断。

\ProseParagraphBreak

最大值原理控制有界经典解，Fourier 能量控制平方可积解；无界空间中的唯一性必须附带相应的解类。热半群在有限指数的 \(L^p\)、无穷远消失的连续函数空间及有界一致连续空间上强连续，在整个 \(L^\infty\) 上则失败。生成元把这些空间中的零时变化记录为带定义域的 Laplace 算子，正时间收缩并不意味着初温反演稳定。
\end{chaptersummary}


\BookExercises

下面八题分别训练级数计算、非齐次边界、镜像法、尺度、长时间质量轮廓、逆问题、各向异性和导数能量。

% 自拟：多项式初温的系数计算，并由已证收敛给出初温与衰减。
\begin{exercise}\label{b2:ex:12-parabolic-profile}
在 \((0,\ell)\) 上取零 Dirichlet 条件及初温 \(f(x)=x(\ell-x)\)。求热解的正弦级数，判断其在零时是否一致取得初值，并求 \(t\to\infty\) 时的第一主项。
\end{exercise}
\begin{solution}
两次分部积分给出非单位正弦系数
\[
 b_n=\dfrac2\ell\int_0^\ell x(\ell-x)\sin\dfrac{n\pi x}\ell\,\mathrm{d}x
 =\dfrac{4\ell^2(1-(-1)^n)}{n^3\pi^3}.
\]
故解为 \(\sum_{n\geq1}b_ne^{-\kappa(n\pi/\ell)^2t}\sin(n\pi x/\ell)\)。系数绝对可和，级数在 \(t\geq0\) 上一致收敛；零时由 \(L^2\) 唯一性识别为连续函数 \(f\)，所以也一致取得初值。第一主项为 \((8\ell^2/\pi^3)e^{-\kappa\pi^2t/\ell^2}\sin(\pi x/\ell)\)，其余只含奇数 \(n\geq3\)，一致余项为 \(O(e^{-9\kappa\pi^2t/\ell^2})\)。
\end{solution}

% 自拟：固定非零边界下的平衡与平均值，区分绝热守恒。
\begin{exercise}\label{b2:ex:12-affine-equilibrium}
端点温度固定为 \(a,b\in\mathbb R\)，初温 \(f\in L^2(0,\ell)\)。求长期平衡温度，并证明解在 \(L^2\) 中指数趋于该平衡。平衡的平均温度是否必须等于初始平均温度？
\end{exercise}
\begin{solution}
令 \(r(x)=a+(b-a)x/\ell\)，则 \(r''=0\) 且满足两端条件。差 \(v=u-r\) 为零边界热解，初值 \(f-r\)，所以
\[
 \|u(t)-r\|_2\leq e^{-\kappa\pi^2t/\ell^2}\|f-r\|_2.
\]
平衡平均为 \((a+b)/2\)，不必等于 \(\ell^{-1}\int f\)。端点可以持续吸收或补给热量，热量变化由法向通量决定；固定端温不要求通量为零。
\end{solution}

% 自拟：半直线核的严格正性与总质量损失，连接镜像和边界吸收。
\begin{exercise}\label{b2:ex:12-half-line}
在 \(x>0\) 上取 \(u(t,0)=0\)。证明相应核为 \(K(t,x,y)=H_{\kappa t}(x-y)-H_{\kappa t}(x+y)\)，并且 \(K>0\) 当 \(x,y,t>0\)。计算 \(\int_0^\infty K(t,x,y)\,\mathrm{d}x\)，解释其严格小于一。
\end{exercise}
\begin{solution}
把初值奇延拓到实线后作热卷积，拆成两个半轴便得到核。因为 \(|x-y|<x+y\)，Gauss 函数随正半径严格递减，所以核严格为正。换元得
\[
 \int_0^\infty K(t,x,y)\,\mathrm{d}x
 =\int_{-y}^{y}H_{\kappa t}(z)\,\mathrm{d}z<1.
\]
单位点热量在正时间有一部分经零温边界流出。奇延拓的负半轴是计算工具；半直线上真正的温度仍为正。
\end{solution}

% 自拟：尺度限制Lp到Lq衰减指数，增加对估计最优尺度的辨析。
\begin{exercise}\label{b2:ex:12-scaling-exponent}
设对全部 \(t>0\)、\(f\in C_c^\infty(\mathbb R^d)\) 有 \(\|S(t)f\|_q\leq Ct^{-\beta}\|f\|_p\)，其中常数与 \(t,f\) 无关。证明必须有 \(\beta=d(1/p-1/q)/2\)。
\end{exercise}
\begin{solution}
固定非零非负 \(f\)，令 \(f_R(x)=f(x/R)\)。热核缩放给出 \(S(R^2)f_R(x)=S(1)f(x/R)\)。代入估计得
\[
 R^{d/q}\|S(1)f\|_q\leq C R^{-2\beta+d/p}\|f\|_p.
\]
两边不随 \(R\) 变化的范数均有限而正。分别令 \(R\downarrow0\)、\(R\to\infty\)，两次要求的幂不等式只能同时以等式满足，故得到所述值。无穷指数按 \(d/\infty=0\) 处理。
\end{solution}

% 自拟：扩散尺度上的质量轮廓，非正文点态衰减的重复。
\begin{exercise}\label{b2:ex:12-mass-asymptotic}
设 \(f\in L^1(\mathbb R^d)\)，\(M=\int f\)。证明
\[
 t^{d/2}u(t,\sqrt t\,\cdot)\longrightarrow M H_\kappa
 \quad\text{于 }L^1(\mathbb R^d),\quad t\to\infty.
\]
若 \(M=0\) 且 \(\int |y|\cdot|f(y)|\,\mathrm{d}y<\infty\)，再证 \(\|u(t)\|_\infty\leq Ct^{-(d+1)/2}\int |y|\cdot|f(y)|\,\mathrm{d}y\)。
\end{exercise}
\begin{solution}
缩放后左侧函数等于 \(\int H_\kappa(z-y/\sqrt t)\cdot f(y)\,\mathrm{d}y\)。与 \(MH_\kappa\) 相减，其 \(L^1\) 范数至多为
\[
 \int\|H_\kappa(\,\cdot-y/\sqrt t)-H_\kappa\|_1\cdot|f(y)|\,\mathrm{d}y.
\]
对固定 \(y\)，平移连续性使范数趋零；它始终不超过二，由控制收敛得第一式。若 \(M=0\)，在原卷积中减去 \(H_{\kappa t}(x)M\)，由微积分基本定理
\[
 |H_{\kappa t}(x-y)-H_{\kappa t}(x)|
 \leq |y|\cdot\|\nabla H_{\kappa t}\|_\infty
 \leq Ct^{-(d+1)/2}|y|.
\]
积分即得改进的衰减。
\end{solution}

% 自拟：正时间算子单射而逆不连续，用归一化高频数据作证。
\begin{exercise}\label{b2:ex:12-backward-instability}
对固定 \(T>0\)，证明 \(S(T)\) 在 \(L^2(\mathbb R^d)\) 上单射，并构造 \(\|f_n\|_2=1\) 而 \(\|S(T)f_n\|_2\to0\) 的序列。说明为何这否定了从终温到初温的有界逆。
\end{exercise}
\begin{solution}
Fourier 乘子 \(e^{-4\pi^2\kappa T|\xi|^2}\) 处处严格为正，所以 \(S(T)f=0\) 迫使 \(\widehat f=0\)。选 \(\widehat f_n\) 支于球 \(B(ne_1,1/2)\)，并归一化其 \(L^2\) 范数为一。则
\[
 \|S(T)f_n\|_2\leq e^{-4\pi^2\kappa T(n-1/2)^2}\longrightarrow0.
\]
若在值域上有有界逆，就有 \(1=\|f_n\|_2\leq C\|S(T)f_n\|_2\)，矛盾。单射保证精确数据至多对应一个初温，但不控制微小终温误差的放大。
\end{solution}

% 自拟：两种正定矩阵不必同时对角化，频率指数相加即可。
\begin{exercise}\label{b2:ex:12-anisotropic-convolution}
设 \(B,C\) 为实对称正定矩阵。证明 \(H_s^B*H_t^C=H_1^{sB+tC}\)，即使 \(BC\ne CB\) 也成立。由此求扩散矩阵为 \(B\)、初值为 \(H_1^C\) 的解。
\end{exercise}
\begin{solution}
两个核的变换分别是 \(e^{-4\pi^2s\xi^TB\xi}\)、\(e^{-4\pi^2t\xi^TC\xi}\)，相乘为 \(e^{-4\pi^2\xi^T(sB+tC)\xi}\)。矩阵 \(sB+tC\) 仍正定，Fourier 唯一性给出核等式；这里相加的是标量二次型，不需要矩阵交换。所求解为 \(H_1^{C+tB}\)，其椭球轴可以随时间旋转，不一定保持初值的主轴。
\end{solution}

% 自拟：对已有能量增加一阶导数，核对总耗散的精确常数。
\begin{exercise}\label{b2:ex:12-gradient-dissipation}
若 \(f\in H^1(\mathbb R^d)\)，证明
\[
 \|\nabla u(T)\|_2^2
 +2\kappa\int_0^T\|\Delta u(t)\|_2^2\,\mathrm{d}t
 =\|\nabla f\|_2^2.
\]
进而计算 \(\int_0^\infty\|\Delta u(t)\|_2^2\,\mathrm{d}t\)。
\end{exercise}
\begin{solution}
对能量证明中的逐频率恒等式乘以 \(4\pi^2|\xi|^2\)，非负换序和 Plancherel 给出所述式子。\(H^1\) 条件保证右端有限，所以不需要预设初值属于 \(H^2\)。当 \(T\to\infty\)，梯度平方范数由频率控制收敛趋零，故积分等于 \(\|\nabla f\|_2^2/(2\kappa)\)。零频点的测度为零，不留下额外项。
\end{solution}
